\documentclass[12pt]{article}
\usepackage[utf8]{inputenc}
\usepackage[T1]{fontenc}
\usepackage{hyperref}
\title{Kobuka Language Resources for K-12 Curriculum Development}
\author{Compiled by Claude Sonnet 5 (claude-sonnet-5) for Andr\'as Kornai}
\date{}
\begin{document}
\maketitle

\section{Introduction}
Kobuka (ISO 639-3 urg; also listed as Kopuka and Wakoro) is spoken in Madang Province, Papua New Guinea, by about 1,400 people. The search results report two varieties, Undi and Wakaro. The ASJP database places it in the Rai Coast group of Trans-New Guinea. Take care not to confuse it with Kobon (kpw) or Kobol (kgu), which are different languages. Very little is documented about Kobuka.

\section{Primary Educational Resources}
No primers, textbooks or school materials in Kobuka turned up in the search. None is known to exist.

\section{Linguistic and Reference Materials}
\begin{itemize}
\item ASJP wordlist entry for Kobuka. It is based on an unpublished manuscript wordlist (Ramu Wordlists, compiled by M.\ S.\ Dryer): \url{https://asjp.clld.org/languages/KOBUKA}
\end{itemize}
No dictionary or grammar was found.

\section{Religious and Community Resources}
None found.

\section{Digital and Audio Resources}
\begin{itemize}
\item Global Recordings Network language page for urg. It gives the name variants, the population figure and the two dialect names, and reports no audio recordings in the language: \url{https://globalrecordings.net/en/language/urg}
\end{itemize}

\section{Curriculum Development Context}
Kobuka appears to be almost undocumented. A curriculum would first need community-led documentation: a wordlist, an orthography and recorded texts. Partnership with PNG education authorities and linguists would be needed. This report cannot name any such partners.
\end{document}